\documentclass[10pt,a4paper]{amsart}
\usepackage{amsmath,amssymb,amsthm,graphicx}
\usepackage{verbatim}
\newenvironment{reference}{\comment}{\endcomment}
%\newenvironment{reference}{}{}
\newenvironment{slogan}{\comment}{\endcomment}
\newenvironment{history}{\comment}{\endcomment}

% For commutative diagrams we use Xy-pic
\usepackage[all]{xy}

% We use 2cell for 2-commutative diagrams.
\xyoption{2cell}
% Theorem environments.
%
\theoremstyle{plain}
\newtheorem{theorem}[subsection]{Theorem}
\newtheorem{proposition}[subsection]{Proposition}
\newtheorem{lemma}[subsection]{Lemma}

\theoremstyle{definition}
\newtheorem{definition}[subsection]{Definition}
\newtheorem{example}[subsection]{Example}
\newtheorem{exercise}[subsection]{Exercise}
\newtheorem{situation}[subsection]{Situation}

\theoremstyle{remark}
\newtheorem{remark}[subsection]{Remark}
\newtheorem{remarks}[subsection]{Remarks}

\numberwithin{equation}{subsection}

% Macros
%
\def\lim{\mathop{\mathrm{lim}}\nolimits}
\def\colim{\mathop{\mathrm{colim}}\nolimits}
\def\Spec{\mathop{\mathrm{Spec}}}
\def\Hom{\mathop{\mathrm{Hom}}\nolimits}
\def\Ext{\mathop{\mathrm{Ext}}\nolimits}
\def\SheafHom{\mathop{\mathcal{H}\!\mathit{om}}\nolimits}
\def\SheafExt{\mathop{\mathcal{E}\!\mathit{xt}}\nolimits}
\def\Sch{\mathit{Sch}}
\def\Mor{\mathop{\mathrm{Mor}}\nolimits}
\def\Ob{\mathop{\mathrm{Ob}}\nolimits}
\def\Sh{\mathop{\mathit{Sh}}\nolimits}
\def\NL{\mathop{N\!L}\nolimits}
\def\CH{\mathop{\mathrm{CH}}\nolimits}
\def\proetale{{pro\text{-}\acute{e}tale}}
\def\etale{{\acute{e}tale}}
\def\QCoh{\mathit{QCoh}}
\def\Ker{\mathop{\mathrm{Ker}}}
\def\Im{\mathop{\mathrm{Im}}}
\def\Coker{\mathop{\mathrm{Coker}}}
\def\Coim{\mathop{\mathrm{Coim}}}

% Boxtimes
%
\DeclareMathSymbol{\boxtimes}{\mathbin}{AMSa}{"02}

%
% Macros for moduli stacks/spaces
%
\def\QCohstack{\mathcal{QC}\!\mathit{oh}}
\def\Cohstack{\mathcal{C}\!\mathit{oh}}
\def\Spacesstack{\mathcal{S}\!\mathit{paces}}
\def\Quotfunctor{\mathrm{Quot}}
\def\Hilbfunctor{\mathrm{Hilb}}
\def\Curvesstack{\mathcal{C}\!\mathit{urves}}
\def\Polarizedstack{\mathcal{P}\!\mathit{olarized}}
\def\Complexesstack{\mathcal{C}\!\mathit{omplexes}}
% \Pic is the operator that assigns to X its picard group, usage \Pic(X)
% \Picardstack_{X/B} denotes the Picard stack of X over B
% \Picardfunctor_{X/B} denotes the Picard functor of X over B
\def\Pic{\mathop{\mathrm{Pic}}\nolimits}
\def\Picardstack{\mathcal{P}\!\mathit{ic}}
\def\Picardfunctor{\mathrm{Pic}}
\def\Deformationcategory{\mathcal{D}\!\mathit{ef}}
\usepackage{cleveref}
\crefname{originallemma}{original lemma}{original lemmas}
\Crefname{originallemma}{Original lemma}{Original lemmas}
\makeatletter
\newlabel{lemma-NAK}{{00DV}{source}{}{sourcepointer-1}{}}
\newlabel{lemma-NAK@cref}{{[originallemma][0][]00DV}{[1][0][]source}{}{}{}}
\makeatother
\crefname{originallemma}{original lemma}{original lemmas}
\Crefname{originallemma}{Original lemma}{Original lemmas}
\makeatletter
\newlabel{lemma-ff}{{00HP}{source}{}{sourcepointer-2}{}}
\newlabel{lemma-ff@cref}{{[originallemma][0][]00HP}{[1][0][]source}{}{}{}}
\makeatother
\crefname{originallemma}{original lemma}{original lemmas}
\Crefname{originallemma}{Original lemma}{Original lemmas}
\makeatletter
\newlabel{lemma-flat-localization}{{00HT}{source}{}{sourcepointer-3}{}}
\newlabel{lemma-flat-localization@cref}{{[originallemma][0][]00HT}{[1][0][]source}{}{}{}}
\makeatother
\title{Generalisations lift along the support of a finite flat module}
\author{The Stacks Project Authors}
\date{stacks-project:algebra.tex:local-snapshot-20261012; GNU FDL 1.2 or later}
\begin{document}
\maketitle
\noindent\textit{Editorial note.} All delivered mathematical text is the author's own, copied from the pinned Stacks Project chapter file algebra.tex; no statement or proof step was rewritten, reordered or removed.\par\medskip
\noindent\textit{Editorial note.} Preamble provenance, as required for this source: the frozen snapshot stages the chapter's own preamble file beside it (preamble.tex, 217 lines, 8116 bytes, the same shared Stacks-wide preamble the published chow and etale-cohomology cards use); the first line of the pinned chapter file is \textbackslash{}input\{preamble\}, so no preamble text lives inside the chapter itself. The card takes that file's lines 8-18 (the verbatim package and the three \textbackslash{}comment-based environments reference/slogan/history at lines 9-12, section 13.2, together with the Xy-pic setup) and its lines 154 to the runtime end of the file (the native theorem declarations and mathematical macros). Not one character was changed, no package was added, no macro was defined, and the excerpt serves only as preamble, never as a span coordinate.\par\medskip
\noindent\textit{Editorial note.} The same-chapter pointers are located at their \textbackslash{}label lines inside the pinned chapter file; the cross-chapter pointers are located in the project's own tag list staged in the snapshot (sources/stacks-chow/tags.txt), at the verbatim lines carrying their canonical tags.\par\medskip
\noindent\textit{Editorial note.} The counted proof refers to same-chapter statements by label; they are recorded as navigation stubs pointed at their \textbackslash{}label lines, and the arguments consuming them stay inside the counted span.\par\medskip
\noindent\textit{Editorial note.} The counted spans contain no \textbackslash{}cite, so the card carries no bibliography.\par\medskip
\noindent\textit{Allowed standard prerequisites (editorial).}\begin{itemize}
\item Flatness as preservation of exactness. Used in the counted proof (lines 10080-10095). Keeps the exact sequences exact after tensoring, which is what produces the required prime.
\item Support of a finite module over the base. Used to translate the topological statement (lines 10076-10080). Identifies the primes in the support with the nonvanishing localisations.
\end{itemize}\medskip
\ifdefined\hypertarget\hypertarget{sourcepointer-1}{}\fi
\noindent\textit{Source reference.} Original lemma 00DV: sources/stacks-algebra/pinned-original-source.tex, source lines 3833--3833. Navigation pointer for \textbackslash{}ref\{lemma-NAK\} used inside the counted span (line 10086); the argument that consumes it is retained in the span.\par\medskip
\ifdefined\hypertarget\hypertarget{sourcepointer-2}{}\fi
\noindent\textit{Source reference.} Original lemma 00HP: sources/stacks-algebra/pinned-original-source.tex, source lines 9341--9341. Navigation pointer for \textbackslash{}ref\{lemma-ff\} used inside the counted span (line 10088); the argument that consumes it is retained in the span.\par\medskip
\ifdefined\hypertarget\hypertarget{sourcepointer-3}{}\fi
\noindent\textit{Source reference.} Original lemma 00HT: sources/stacks-algebra/pinned-original-source.tex, source lines 9413--9413. Navigation pointer for \textbackslash{}ref\{lemma-flat-localization\} used inside the counted span (line 10082); the argument that consumes it is retained in the span.\par\medskip
\begin{lemma}
\label{lemma-going-down-flat-module}
Let $R \to S$ be a ring map. Let $N$ be a finite $S$-module flat over $R$.
Endow $\text{Supp}(N) \subset \Spec(S)$ with the induced topology.
Then generalizations lift along $\text{Supp}(N) \to \Spec(R)$.
\end{lemma}

\begin{proof}
The meaning of the statement is as follows. Let
$\mathfrak p \subset \mathfrak p' \subset R$ be primes. Let
$\mathfrak q' \subset S$ be a prime $\mathfrak q' \in \text{Supp}(N)$
Then there exists a prime $\mathfrak q \subset \mathfrak q'$,
$\mathfrak q \in \text{Supp}(N)$ lying over $\mathfrak p$.
As $N$ is flat over $R$ we see that $N_{\mathfrak q'}$ is flat
over $R_{\mathfrak p'}$, see Lemma \ref{lemma-flat-localization}.
As $N_{\mathfrak q'}$ is finite over $S_{\mathfrak q'}$
and not zero since $\mathfrak q' \in \text{Supp}(N)$ we see
that $N_{\mathfrak q'} \otimes_{S_{\mathfrak q'}} \kappa(\mathfrak q')$
is nonzero by Nakayama's Lemma \ref{lemma-NAK}.
Thus $N_{\mathfrak q'} \otimes_{R_{\mathfrak p'}} \kappa(\mathfrak p')$
is also not zero. We conclude from Lemma \ref{lemma-ff}
that $N_{\mathfrak q'} \otimes_{R_{\mathfrak p'}} \kappa(\mathfrak p)$
is nonzero. Let
$J \subset S_{\mathfrak q'} \otimes_{R_{\mathfrak p'}} \kappa(\mathfrak p)$
be the annihilator of the finite nonzero module
$N_{\mathfrak q'} \otimes_{R_{\mathfrak p'}} \kappa(\mathfrak p)$.
Since $J$ is a proper ideal we can choose a prime $\mathfrak q \subset S$
which corresponds to a prime of
$S_{\mathfrak q'} \otimes_{R_{\mathfrak p'}} \kappa(\mathfrak p)/J$.
This prime is in the support of $N$, lies over $\mathfrak p$, and
is contained in $\mathfrak q'$ as desired.
\end{proof}

\end{document}
